\documentclass{article}

\usepackage{listings}
\usepackage{xcolor}
\renewcommand{\ttdefault}{pcr}
\usepackage{inconsolata}
%\usepackage{DejaVuSansMono}   % ← replaces Computer Modern Typewriter with a
                            %   modern, clean monospace font for all \texttt{}
                            %   and lstlisting environments automatically

% Define the colors to match the paper's style
\definecolor{codekeyword}{RGB}{30, 100, 200}      % blue for keywords like for, int
\definecolor{codecomment}{RGB}{0, 128, 96}         % teal/green for // comments
\definecolor{codepragma}{RGB}{150, 50, 150}        % purple for #pragma directives
\definecolor{codebackground}{RGB}{248, 248, 248}   % very light gray background
\definecolor{linehighlight}{RGB}{230, 230, 230}    % gray for highlighted lines

% Configure the global listing style — this applies to ALL lstlisting blocks
\lstset{
  language=C++,
  basicstyle=\ttfamily\small,          % monospaced font, small size
  keywordstyle=\color{codekeyword}\bfseries,   % blue + bold for C++ keywords
  commentstyle=\color{codecomment}\itshape,    % teal italic for comments
  stringstyle=\color{red!70!black},            % dark red for string literals
  % Add HLS-specific keywords so they get colored too
  morekeywords={pragma, HLS, PIPELINE, UNROLL, ARRAY_PARTITION,
                II, factor, cyclic, complete, dataflow},
  numbers=left,                        % line numbers on the left margin
  numberstyle=\tiny\color{gray},       % line numbers are small and gray
  stepnumber=1,                        % number every single line
  numbersep=8pt,                       % space between line numbers and code
  backgroundcolor=\color{codebackground},  % light gray background for code block
  frame=single,                        % draw a box border around the code
  rulecolor=\color{gray!40},           % border color
  showstringspaces=false,              % don't show space markers inside strings
  breaklines=true,                     % wrap long lines instead of overflowing
  tabsize=2,                           % treat tabs as 2 spaces
  captionpos=b,                        % put the caption below the listing
}

\begin{document}

\title{Test Document}
\author{Aaron}
\date{\today}
\maketitle

Hello, this is a test. Here is some inline math: $E = mc^2$.

\begin{equation}
  \sum_{i=0}^{N} x_i \cdot w_i
\end{equation}

% The label= option lets you reference this listing elsewhere with \ref{}
\begin{lstlisting}[
  caption={Pipelined dot product loop with HLS pragma.},
  label={lst:dot_product}
]
for (int i = 0; i < N; i++) {
    #pragma HLS PIPELINE II=1
    result += x[i] * w[i];  // accumulate weighted sum
}
\end{lstlisting}

\end{document}\documentclass{article}

\usepackage{listings}
\usepackage{xcolor}

% Define the colors to match the paper's style
\definecolor{codekeyword}{RGB}{30, 100, 200}      % blue for keywords like for, int
\definecolor{codecomment}{RGB}{0, 128, 96}         % teal/green for // comments
\definecolor{codepragma}{RGB}{150, 50, 150}        % purple for #pragma directives
\definecolor{codebackground}{RGB}{248, 248, 248}   % very light gray background
\definecolor{linehighlight}{RGB}{230, 230, 230}    % gray for highlighted lines

% Configure the global listing style — this applies to ALL lstlisting blocks
\lstset{
  language=C++,
  basicstyle=\ttfamily\small,          % monospaced font, small size
  keywordstyle=\color{codekeyword}\bfseries,   % blue + bold for C++ keywords
  commentstyle=\color{codecomment}\itshape,    % teal italic for comments
  stringstyle=\color{red!70!black},            % dark red for string literals
  % Add HLS-specific keywords so they get colored too
  morekeywords={pragma, HLS, PIPELINE, UNROLL, ARRAY_PARTITION,
                II, factor, cyclic, complete, dataflow},
  numbers=left,                        % line numbers on the left margin
  numberstyle=\tiny\color{gray},       % line numbers are small and gray
  stepnumber=1,                        % number every single line
  numbersep=8pt,                       % space between line numbers and code
  backgroundcolor=\color{codebackground},  % light gray background for code block
  frame=single,                        % draw a box border around the code
  rulecolor=\color{gray!40},           % border color
  showstringspaces=false,              % don't show space markers inside strings
  breaklines=true,                     % wrap long lines instead of overflowing
  tabsize=2,                           % treat tabs as 2 spaces
  captionpos=b,                        % put the caption below the listing
}

\begin{document}

\title{Test Document}
\author{Aaron}
\date{\today}
\maketitle

Hello, this is a test. Here is some inline math: $E = mc^2$.

\begin{equation}
  \sum_{i=0}^{N} x_i \cdot w_i
\end{equation}

% The label= option lets you reference this listing elsewhere with \ref{}
\begin{lstlisting}[
  caption={Pipelined dot product loop with HLS pragma.},
  label={lst:dot_product}
]
for (int i = 0; i < N; i++) {
    #pragma HLS PIPELINE II=1
    result += x[i] * w[i];  // accumulate weighted sum
}
\end{lstlisting}

\end{document}
